\section{引言}\label{sec:intro}

核子的部分子分布函数（PDF）不仅是刻画强子内部结构的重要物理量，也是对高能散射过程作出预言的关键要素~\cite{Butterworth:2015oua,Alekhin:2017kpj,Gao:2017yyd}。因此，从第一性原理计算 PDF 一直是核物理与粒子物理中的“圣杯”。由于 PDF 内嵌于强子中低能的夸克与胶子自由度，它们涉及强相互作用的红外（IR）动力学，只能由格点 QCD 这类非微扰方法确定。

在 QCD 因子化~\cite{Collins:2011zzd}框架下，夸克 PDF 定义为
\begin{align} \label{eq:lcpdf}
q(x,\mu) \equiv\! \int\! {d\xi^-\over 4\pi} e^{-ixP^+\xi^-}\! \langle P | \bar{\psi}  (\xi^-) \gamma^+ U (\xi^-, 0) \psi (0) | P \rangle ,
\end{align}
其中 $|P\rangle$ 表示动量为 $P_\mu=(P_t,0,0,P_z)$ 的核子态。$x$ 是夸克动量份额，$\mu$ 是 $\overline{\rm MS}$ 方案中的重正化标度。$\xi^{\pm}=(t\pm z)/\sqrt{2}$ 为光锥坐标。为保持规范不变性，引入了类光 Wilson 线：
\begin{align}
U(\xi^-,0) = P\exp\biggl(-ig\int_0^{\xi^-} d\eta^- A^+(\eta^-)\biggr)\ .
\end{align}

PDF 由光锥坐标定义，但格点模拟只能在欧氏空间中进行，无法恰当处理涉及真实时间的光锥量。因此，在欧氏格点上模拟 PDF 是一项极其困难的任务。早期基于算符乘积展开（OPE）的研究只能给出 PDF 的最低几个矩~\cite{Martinelli:1987zd,Martinelli:1988xs,Detmold:2001dv,Dolgov:2002zm}。

最近，文献~\cite{Ji:2013dva}提出了一种可以直接从格点 QCD 获得 PDF 的 $x$ 依赖关系的新方法，现在被表述为大动量有效理论（LaMET）~\cite{Ji:2014gla}。在该框架下，人们可以从核子态中某些静态算符的关联函数出发提取 PDF 以及其他光锥物理量。一方面，这些静态关联（通常称为准观测量，quasi observables）可以直接在欧氏格点上计算，并动态地依赖于核子动量。另一方面，在大动量下，准观测量可以因子化为部分子观测量与一个微扰匹配系数的卷积，修正被大核子动量的幂次压低。将两侧的结果等同起来，就提供了一条确定光锥 PDF 的直接途径。

要在 LaMET 中计算夸克 PDF，需从“准 PDF”出发，其定义为夸克沿 $z$ 方向的等时关联函数~\cite{Ji:2013dva}：
\begin{align}\label{eq:qpdf}
\widetilde{q}_{\Gamma}(x,P_z) \equiv \int_{-\infty}^\infty {dz\over 4\pi} e^{ixP_zz}\langle P |O_{\Gamma}(z)| P \rangle \,.
\end{align}
上式中，$O_{\Gamma}(z)=\bar{\psi}  (z) \Gamma U (z, 0) \psi (0)$，$\Gamma=\gamma^z$ 或 $\Gamma=\gamma^t$，类空 Wilson 线为
\begin{align}
U(z,0) = P\exp\left(-ig\int_0^z dz' A_z(z')\right) .
\end{align}
对有限但很大的动量 $P_z$，$\widetilde{q}(x,P_z)$ 的支撑区间为 $-\infty < x < \infty$。与具有 boost 不变性的光锥 PDF 不同，准 PDF 对核子动量 $P_z$ 有非平庸的依赖。在以正规化无关动量减除（RI/MOM）等方案将准 PDF 重正化之后，可以通过因子化定理把重正化的准 PDF 匹配到 $\overline{\rm MS}$ PDF~\cite{Ji:2013dva,Ji:2014gla,Stewart:2017tvs,Izubuchi:2018srq,Ma:2014jla,Ma:2017pxb}：
\begin{align} \label{eq:fact}
\widetilde{q}(x,P_z, p^R_z,\mu_R)=&\int_{-1}^1 {dy\over |y|}\: C\left({x\over y},r,\frac{yP_z}{\mu},\frac{yP_z}{p_z^R}\right) \, q(y,\mu)\nonumber\\
&+\mathcal{O}\left({M^2\over P_z^2},{\Lambda_{\text{QCD}}^2\over x^2 P_z^2}\right) ,
\end{align}
其中 RI/MOM 方案引入了 $p_z^R$ 与 $\mu_R$：$p_z^R$ 是所涉及部分子的动量，$\mu_R$ 是重正化标度。$r=\mu_R^2/(p_z^R)^2$；$C$ 为微扰匹配系数；$\mathcal{O}(M^2/P_z^2, \Lambda_{\rm QCD}^2/x^2P_z^2)$ 表示被大核子动量幂次压低的核子质量贡献与更高扭度贡献。$q$、$\widetilde{q}$ 与 $C$ 中的味指标均已隐含。当 $-1<y<0$ 时，分布指反夸克分布。

自 LaMET 提出以来，理论与格点计算两方面都取得了显著进展。应当指出，这些发展是以相互促进的方式实现的。LaMET 首先被用于计算质子同位旋矢量夸克分布 $f_{u-d}$~\cite{Lin:2014zya,Alexandrou:2015rja,Chen:2016utp,Alexandrou:2016jqi,Chen:2018fwa,Alexandrou:2018eet}，包括非极化、极化与横向度情形，随后又推广到介子分布振幅~\cite{Zhang:2017bzy,Chen:2017gck}。最早的格点研究采用的是横动量截断方案下一圈阶的匹配系数~\cite{Xiong:2013bka,Ji:2015qla,Xiong:2015nua}。然而，正如文献~\cite{Xiong:2013bka}所指出的，最初的准 PDF 存在紫外（UV）线性发散，这可能给其格点矩阵元素的重正化带来严重问题~\cite{Li:2016amo,Rossi:2017muf,Rossi:2018zkn}。此后人们的注意力集中到其重正化性质上~\cite{Ji:2015jwa,Ishikawa:2016znu,Chen:2016fxx,Xiong:2017jtn,Constantinou:2017sej,Ji:2017oey,Ishikawa:2017faj,Green:2017xeu,Spanoudes:2018zya}，最终证明坐标空间中准 PDF 在连续极限下对强耦合常数 $\alpha_s$ 的所有阶都具有可乘重正化性~\cite{Ji:2017oey,Ishikawa:2017faj}。这一发现进一步推动了准 PDF 非微扰重正化（NPR）的格点分析~\cite{Alexandrou:2017huk,Chen:2017mzz,Green:2017xeu}（采用 RI/MOM 方案~\cite{Martinelli:1994ty}），以及 RI/MOM 准 PDF 与 $\overline{\rm MS}$ PDF 之间匹配系数的计算~\cite{Stewart:2017tvs}。除重正化外，有限核子质量修正也被算到 $M^2/P_z^2$ 的所有阶~\cite{Chen:2016utp}，而更高扭度的 $\mathcal O(\Lambda^2_{\rm QCD}/x^2P_z^2)$ 效应则通过把若干 $P_z$ 值处的结果外推到无穷大动量而被数值地消除~\cite{Lin:2014zya,Chen:2016utp}。基于这些研究，物理 pion 质量下的同位旋矢量夸克 PDF 计算已成为可能~\cite{Lin:2017ani,Alexandrou:2018pbm,Chen:2018xof,Alexandrou:2018eet}。人们在准 PDF 格点重正化中可能的算符混合也做了研究~\cite{Constantinou:2017sej,Alexandrou:2017huk,Green:2017xeu,Chen:2017mzz}，其混合模式已在文献~\cite{Chen:2017mie}中被分类。针对长距离处空间关联函数傅里叶变换的系统误差控制方法见文献~\cite{Lin:2017ani,Chen:2017lnm}。LaMET 也已被尝试用于研究横动量依赖分布（TMD）~\cite{Ji:2014hxa,Ji:2018hvs,Ebert:2018gzl,Ebert:2019okf,Ebert:2019tvc,Ji:2019sxk,Shanahan:2019zcq,Ji:2019ewn}以及胶子 PDF~\cite{Wang:2017qyg,Wang:2017eel,Fan:2018dxu,Zhang:2018diq,Li:2018tpe,Wang:2019tgg}。

除 LaMET 外，近年来还提出了其他一些从格点 QCD 计算 PDF 的有趣方法。例如，可以从一类“格点截面”（lattice cross sections）提取 PDF~\cite{Ma:2014jla,Ma:2017pxb}；梯度流方法中的涂抹准 PDF 则被提出来扫除格点计算中的幂次发散~\cite{Monahan:2016bvm,Monahan:2017hpu}。还可以研究赝分布（pseudo distribution）~\cite{Radyushkin:2017cyf}，它通过傅里叶变换与准 PDF 相联系。该方法展现出有趣的重正化性质~\cite{Orginos:2017kos,Karpie:2017bzm}，而就到 PDF 的因子化而言它与 LaMET 一致~\cite{Ji:2017rah,Zhang:2018ggy,Izubuchi:2018srq}。此外，还有利用流—流关联函数计算强子张量~\cite{Liu:1993cv,Liang:2017mye}，以及计算 PDF 高阶矩、光锥分布振幅等的建议~\cite{Detmold:2005gg,Braun:2007wv,Liang:2017mye,Chambers:2017dov,Bali:2018spj,Bali:2019ecy}。这些不同方法各有自身的系统误差，但彼此之间可以互相比较。

曾有观点认为定域矩算符之间的幂次发散混合可能破坏准 PDF 的重正化~\cite{Rossi:2017muf,Rossi:2018zkn}，但在 LaMET 中该问题并不存在，因为人们只需先对格点上重正化后的准 PDF 取连续极限，再将其匹配得到 PDF 的 $x$ 依赖关系。因子化已在连续理论中得到严格推导~\cite{Ma:2014jla,Izubuchi:2018srq}，因而只需关注非定域空间关联函数本身的重正化。因此，定域矩算符的重正化与准 PDF 无关。此外，围绕准 PDF 的闵氏与欧氏矩阵元素之间的 LaMET 匹配也曾存在困惑~\cite{Carlson:2017gpk},相关澄清见文献~\cite{Ji:2017rah,Briceno:2017cpo}。

大多数现有的格点计算在非极化准 PDF 中使用 $\Gamma=\gamma^z$（除文献~\cite{Chen:2018xof,Alexandrou:2018pbm} 外），而现在已经知道它会与标量准 PDF 算符 $O_{I}$ 在 $O(a^0)$ 阶发生混合~\cite{Constantinou:2017sej,Chen:2017mie,Green:2017xeu}。这种算符混合给非微扰重正化带来了额外的系统不确定性~\cite{Alexandrou:2017huk,Chen:2017mzz,Green:2017xeu,Lin:2017ani}，从而限制了所提取 PDF 的精度。相反，$\Gamma=\gamma^t$ 的情形在 $O(a^0)$ 阶不与 $O_{I}$ 发生算符混合~\cite{Constantinou:2017sej,Chen:2017mie,Green:2017xeu}。因此，从 $\Gamma=\gamma^t$ 的准 PDF 出发是非常值得的。这正是本研究的主要动机之一。

在本工作中，我们将从 $\Gamma=\gamma^t$ 的准 PDF 出发进行非极化同位旋矢量夸克分布的格点计算，非微扰重正化流程与文献~\cite{Chen:2017mzz}中 $\Gamma=\gamma^z$ 情形相同。计算采用 clover 费米子，规范组态为开边界条件下 $N_f=2+1$（简并的上/下夸克及奇异夸克）味的 CLS 系综，pion 质量 $M_{\pi}=356$~MeV、格距 $a=0.086$~fm~\cite{Bruno:2014jqa}。我们将检验对核子动量 $P_z$ 及 RI/MOM 标度 $p_z^R$、$\mu_R$ 的依赖，以及对 RI/MOM 重正化中砍腿格林函数投影算符选取的依赖。由于误差较大，我们难以分辨早期未经格点重正化的研究中所观察到的海夸克不对称~\cite{Lin:2014zya,Alexandrou:2015rja,Chen:2016utp,Alexandrou:2016jqi}。我们计划在未来分析精度更高的 CLS 系综，直至物理质量和 $a<0.04$~fm，既包括固定 $(m_s+m_u+m_d)$ 于物理值的情形，也包括用味 SU(3) 与 SU(2) 外推取物理 $m_s$ 的情形。

本文其余部分安排如下：第~\ref{sec:matching} 节简要回顾 RI/MOM 方案下准 PDF 的非微扰重正化与匹配流程，特别是 $\Gamma=\gamma^t$ 情形的一圈匹配系数；第~\ref{sec:numerical} 节描述强子矩阵元素及其非微扰重正化的格点模拟细节，并在该节讨论系统误差；第~\ref{sec:final_results} 节给出带有统计与系统不确定度的非极化同位旋矢量夸克 PDF 的 $x$ 依赖结果；最后一节总结全文。
